\documentclass[11pt,letterpaper]{article}

\usepackage[T1]{fontenc}
\usepackage{amsmath,amssymb}
\usepackage{tikz}
\usepackage{pgfplots}
\pgfplotsset{compat=1.18}
\usepackage{algorithm}
\usepackage{algpseudocode}
\usepackage{booktabs}
\usepackage{enumitem}
\usepackage[hidelinks]{hyperref}

\newcommand{\N}{\mathbb{N}}
\newcommand{\Z}{\mathbb{Z}}
\newcommand{\Q}{\mathbb{Q}}
\newcommand{\R}{\mathbb{R}}
\newcommand{\C}{\mathbb{C}}
\newcommand{\Map}[3]{#1 \colon #2 \to #3}
\newcommand{\Mod}[1]{#1\text{-}\mathrm{Mod}}
\newcommand{\GroupHom}{\mathrm{Hom}_{\mathbf{Grp}}}
\newcommand{\RingHom}{\mathrm{Hom}_{\mathbf{Ring}}}

\newtheorem{theorem}{Theorem}
\newtheorem{lemma}[theorem]{Lemma}
\newtheorem{corollary}[theorem]{Corollary}
\newtheorem{definition}[theorem]{Definition}
\newtheorem{remark}[theorem]{Remark}
\newenvironment{proof}{\noindent\textit{Proof.}}{\hfill\(\square\)}

\title{Generalized Collatz Library}
\author{}
\date{\today}

\begin{document}
\maketitle

\begin{abstract}
For every integer \(t\ge0\), the raw Collatz orbits starting at
\(n=11003+1048576t\) and \(p=9915+944784t\) have prefixes that first meet at
\(33466+3188646t\), after \(32\) and \(6\) raw steps respectively.
Their next common value is \(j=16733+1594323t\).
The parent is positive and satisfies \(100p\le91n\).
A finite table of affine coefficients proves the branch legality,
the exact meeting clocks, and the absence of earlier intersections
uniformly in \(t\). These are smaller-parent identities on one cylinder,
not an assertion that its orbits reach \(1\).
\end{abstract}

\section{Raw clock and affine coefficients}

We use the raw map
\[
C(x)=
\begin{cases}
x/2,&x\equiv0\pmod2,\\
3x+1,&x\equiv1\pmod2.
\end{cases}
\]
The map is defined on the nonnegative integers; all states below are positive.
In particular, an odd step and its following division count as two steps.
The shortcut map would count that pair as one.

Fix an integer \(t\ge0\), and write
\[
n=11003+1048576t,\qquad p=9915+944784t,\qquad
j=16733+1594323t.
\]

\begin{lemma}[Affine coefficient rules]
\label{lem:coefficients}
If \(a,b\) are even nonnegative integers, then
\(C(a+bt)=a/2+(b/2)t\).
If \(a\) is odd and \(b\) is even, then
\(C(a+bt)=(3a+1)+(3b)t\).
If \(a<c\) and \(b\le d\), then \(a+bt<c+dt\).
\end{lemma}

\begin{proof}
In the first case, \(a+bt=2(a/2+(b/2)t)\).
In the second, \(bt\) is even, so the argument is odd, and expanding
\(3(a+bt)+1\) gives the stated coefficients.
For the comparison,
\(a+bt\le a+dt<c+dt\), since \(t\ge0\).
\end{proof}

\section{A complete finite coefficient table}

Each label in Table~\ref{tab:states} denotes \(a+bt\) with the displayed
constant \(a\) and slope \(b\). The source labels are \(n_0,\ldots,n_{33}\);
the parent labels are \(p_0,\ldots,p_7\).
The labels \(e,f\) are the two additional forms excluded below.
The table is sorted by its constant: read down the left half, then
down the right half. Its constants strictly increase and its slopes
do not decrease. Thus Lemma~\ref{lem:coefficients} compares any two
distinct rows strictly for every \(t\ge0\).

\begin{table}[ht]
\centering
\small
\begin{tabular}{crrcrr}
\toprule
State & \(a\) & \(b\) & State & \(a\) & \(b\)\\
\midrule
\(e\) & 8251 & 786432 & \(n_{19}\) & 35255 & 3359232 \\
\(p_{0}\) & 9915 & 944784 & \(n_{6}\) & 37138 & 3538944 \\
\(n_{0}\) & 11003 & 1048576 & \(n_{11}\) & 41782 & 3981312 \\
\(p_{5}\) & 11155 & 1062882 & \(p_{3}\) & 44620 & 4251528 \\
\(n_{5}\) & 12379 & 1179648 & \(n_{29}\) & 44621 & 4251528 \\
\(n_{10}\) & 13927 & 1327104 & \(n_{16}\) & 47006 & 4478976 \\
\(p_{2}\) & 14873 & 1417176 & \(n_{3}\) & 49516 & 4718592 \\
\(n_{2}\) & 16505 & 1572864 & \(n_{21}\) & 52883 & 5038848 \\
\(n_{33}, p_{7}\) & 16733 & 1594323 & \(n_{8}\) & 55708 & 5308416 \\
\(n_{7}\) & 18569 & 1769472 & \(n_{26}\) & 59494 & 5668704 \\
\(n_{12}\) & 20891 & 1990656 & \(n_{13}\) & 62674 & 5971968 \\
\(p_{4}\) & 22310 & 2125764 & \(n_{31}\) & 66932 & 6377292 \\
\(n_{17}\) & 23503 & 2239488 & \(n_{18}\) & 70510 & 6718464 \\
\(f\) & 24754 & 2359296 & \(n_{23}\) & 79325 & 7558272 \\
\(n_{4}\) & 24758 & 2359296 & \(n_{28}\) & 89242 & 8503056 \\
\(n_{9}\) & 27854 & 2654208 & \(n_{15}\) & 94012 & 8957952 \\
\(p_{1}\) & 29746 & 2834352 & \(n_{20}\) & 105766 & 10077696 \\
\(n_{27}\) & 29747 & 2834352 & \(n_{25}\) & 118988 & 11337408 \\
\(n_{14}\) & 31337 & 2985984 & \(n_{30}\) & 133864 & 12754584 \\
\(n_{1}\) & 33010 & 3145728 & \(n_{22}\) & 158650 & 15116544 \\
\(n_{32}, p_{6}\) & 33466 & 3188646 & \(n_{24}\) & 237976 & 22674816 \\
\bottomrule
\end{tabular}
\caption{All forty distinct chain forms and the two excluded forms.}
\label{tab:states}
\end{table}

\begin{theorem}[Legal successor chains]
\label{thm:chains}
For every integer \(t\ge0\),
\[
n_i=C^i(n)\quad(0\le i\le33),\qquad
p_i=C^i(p)\quad(0\le i\le7).
\]
Their common rows are precisely
\[
n_{32}=p_6=33466+3188646t,\qquad
n_{33}=p_7=j.
\]
\end{theorem}

\begin{proof}
The initial rows are \(n_0=n\) and \(p_0=p\).
All source slopes before \(n_{33}\) and all parent slopes before \(p_7\)
are even, so each input row has the parity of its constant.
For an even input row, halving both coefficients gives the row with the
next index; for an odd input row, the next coefficients are \(3a+1,3b\).
Every required pair appears in the table. For example,
\[
C(n_4)=C(24758+2359296t)=12379+1179648t=n_5,
\]
and
\[
C(p_5)=C(11155+1062882t)=33466+3188646t=p_6.
\]
Applying the coefficient rules to each consecutive pair proves the chains
by induction on the displayed indices.
Finally, the table has exactly two rows carrying both a source label and
a parent label. All other rows are unequal by coefficient dominance.
\end{proof}

\begin{corollary}[Exact first meeting]
\label{cor:first}
The two displayed prefixes have no intersection before their shared tail.
In particular,
\[
C^{32}(n)=C^6(p),\qquad C^{33}(n)=C^7(p)=j.
\]
\end{corollary}

\begin{proof}
Every source row \(n_i\) with \(i<32\) is distinct from every parent row
\(p_k\) with \(0\le k\le7\), by the two-label description in
Theorem~\ref{thm:chains}. Likewise, the parent rows with \(k<6\) are
distinct from all displayed source rows. The two listed shared rows
therefore give the first intersection of these prefixes.
\end{proof}

\section{Parent size and the terminal parity split}

\begin{theorem}[Positive contracted parent]
\label{thm:parent}
For every integer \(t\ge0\),
\[
0<p<n<j,\qquad 100p\le91n.
\]
Moreover,
\[
n-p=1088+103792t,\qquad j-n=5730+545747t.
\]
\end{theorem}

\begin{proof}
The constant of \(p\) is positive. Subtracting the coefficients gives the
two gap identities, whose right-hand sides are positive.
For the contraction,
\[
91n-100p=9773+942016t>0.
\]
Thus the claimed weak bound holds, with strict slack.
\end{proof}

\begin{remark}[The next branch depends on the parameter]
Both coefficients of \(j\) are odd, so \(j\equiv t+1\pmod2\).
Consequently, \(C(j)=j/2\) when \(t\) is odd, whereas
\(C(j)=3j+1\) when \(t\) is even.
There is no single affine continuation in \(t\) on the full cylinder.
\end{remark}

\begin{corollary}[Excluded forms]
No displayed chain state equals
\[
e=8251+786432t,\qquad f=24754+2359296t.
\]
\end{corollary}

\begin{proof}
The two forms have separate rows in Table~\ref{tab:states} and are
strictly separated from every source and parent row by coefficient
dominance. In particular, \(f\) is not the fourth successor
\(n_4=24758+2359296t\), even though their slopes agree.
\end{proof}

\section{Recovery and scope}

If the parent orbit reaches \(1\), then the source orbit also reaches \(1\).
Indeed, the parent prefix \(p_0,\ldots,p_6\) contains no value \(1\);
its smallest constant is \(9915\).
Any visit to \(1\) therefore occurs at a time \(h>6\).
The identity \(C^{32}(n)=C^6(p)\) then gives
\(C^{32+h-6}(n)=C^h(p)=1\).
This argument uses the same parent and the same forward orbit.

The fixed factor \(91/100<1\) makes this parent replacement
well-founded whenever it can be used. It does not ensure that a later
parent lies in this cylinder, supply a rule for all other seeds, or prove
capture of \(p\).
The theorem is therefore a uniform smaller-parent join for one
explicit arithmetic progression, not a global capture theorem.

\section{Formal proofs and exact computation}

The accompanying \texttt{lean/LocalJoin.lean} proves the coefficient rules,
both successor chains, their cross comparisons, the excluded forms,
the four raw iterate identities, and the parent contraction.
The finite comparison conjunctions retain the actual numerical coefficients;
their proofs use the same dominance argument as the table.

The Go program in \texttt{go/LocalJoin.go} checks all coefficient transitions,
all cross comparisons and both exclusions directly.
It also compares numerical raw iterates with every row for parameters
\(0\le t\le10^4\), with explicit input support through \(t=10^6\).
These finite numerical checks are supplementary:
the proofs above and in Lean cover every integer \(t\ge0\).

\end{document}
